\documentclass[11pt]{article}

\usepackage[a4paper,margin=31mm]{geometry}
\usepackage{amsmath,amssymb,amsthm}
\usepackage{microtype}
\usepackage{enumitem}
\usepackage[dvipsnames]{xcolor}
\usepackage[colorlinks=true,linkcolor=MidnightBlue,urlcolor=MidnightBlue,
            citecolor=MidnightBlue]{hyperref}

\newtheorem{theorem}{Theorem}
\newtheorem{lemma}{Lemma}
\newtheorem{proposition}{Proposition}
\theoremstyle{remark}
\newtheorem*{remark}{Remark}

\newcommand{\HH}[2]{H_{#1}^{(#2)}}
\newcommand{\zetaThree}{\zeta(3)}
\newcommand{\lcm}{\operatorname{lcm}}

\title{\textbf{Apéry's Companion in One Harmonic Sum}\\
\large A short route through the irrationality proof for \(\zeta(3)\)}
\author{}
\date{}

\begin{document}
\maketitle

\begin{abstract}
Apéry's two sequences are usually presented asymmetrically: the integer
sequence \(a_n\) has a short binomial sum, whereas its rational companion
\(b_n\) is given by a nested sum.  We explain how the unexpectedly short
formula
\[
 b_n=\sum_{k=0}^n
 \binom nk^2\binom{n+k}{k}^2
 \bigl(2\HH n3-\HH k3\bigr)
\]
makes two parts of Apéry's proof nearly immediate: the convergence
\(b_n/a_n\to\zetaThree\) and the denominator bound
\(d_n^3b_n\in\mathbb Z\), where \(d_n=\lcm(1,\ldots,n)\).
The Apéry recurrence and a one-line discrete Wronskian calculation then
recover the exponentially small linear forms and complete the irrationality
proof.
\end{abstract}

\section{The short formula and the proof map}

Write
\[
 K_{n,k}:=\binom nk^2\binom{n+k}{k}^2,
 \qquad
 \HH m3:=\sum_{j=1}^m\frac1{j^3},
 \qquad \HH 03:=0.
\]
The two Apéry sequences take the parallel form
\begin{equation}\label{eq:ab}
 a_n=\sum_{k=0}^nK_{n,k},
 \qquad
 b_n=\sum_{k=0}^nK_{n,k}
       \bigl(2\HH n3-\HH k3\bigr).
\end{equation}

The closed form does not by itself create the exponential approximation.
Its role is more precise and, for exposition, especially useful:
\begin{enumerate}[label=\textup{(\roman*)}]
 \item the binomial kernel concentrates at \(k\sim n/\sqrt2\), so
       \eqref{eq:ab} makes \(b_n/a_n\to\zetaThree\) transparent;
 \item the harmonic numbers make \(d_n^3b_n\in\mathbb Z\) termwise;
 \item the recurrence and its discrete Wronskian then force the error
       \(a_n\zetaThree-b_n\) to be exponentially small.
\end{enumerate}

We take as our one recurrence input the classical fact that both sequences in
\eqref{eq:ab} satisfy
\begin{equation}\label{eq:recurrence}
 (n+1)^3u_{n+1}
 =
 (34n^3+51n^2+27n+5)u_n-n^3u_{n-1},
\end{equation}
with
\[
 a_0=1,\quad a_1=5,
 \qquad
 b_0=0,\quad b_1=6.
\]
Thus the only identity not proved in this short note is the recurrence itself;
it may be established independently, for example by creative telescoping.

\section{Where the binomial kernel lives}

For \(0<x<1\), Stirling's formula gives
\begin{equation}\label{eq:phi-limit}
 \frac1n\log K_{n,\lfloor xn\rfloor}\longrightarrow\Phi(x),
\end{equation}
where
\begin{equation}\label{eq:phi}
 \Phi(x)
 =
 2\bigl((1+x)\log(1+x)-2x\log x-(1-x)\log(1-x)\bigr).
\end{equation}
Its derivative is
\[
 \Phi'(x)=2\log\left(\frac{1-x^2}{x^2}\right).
\]
Consequently \(\Phi\) has a unique maximum at
\[
 x_*=\frac1{\sqrt2}.
\]
At this point
\[
 \exp\bigl(\Phi(x_*)\bigr)
 =(1+\sqrt2)^4
 =17+12\sqrt2.
\]
We abbreviate
\begin{equation}\label{eq:lambda}
 \lambda:=17+12\sqrt2.
\end{equation}

The elementary discrete Laplace principle now gives two facts that we shall
use.  First,
\begin{equation}\label{eq:a-growth}
 \lim_{n\to\infty}a_n^{1/n}=\lambda.
\end{equation}
Second, if
\[
 p_{n,k}:=\frac{K_{n,k}}{a_n},
\]
then \(p_{n,k}\) is a probability distribution which concentrates
exponentially near \(k=n/\sqrt2\).  More explicitly, for every
\(\delta>0\),
\begin{equation}\label{eq:concentration}
 \sum_{\left|k/n-1/\sqrt2\right|\geq\delta}p_{n,k}
 \longrightarrow0.
\end{equation}
Indeed, on the compact set separated by \(\delta\) from \(x_*\), the
continuous function \(\Phi\) is smaller than \(\Phi(x_*)\) by a fixed positive
amount.  Equation \eqref{eq:phi-limit}, together with the fact that there are
only \(n+1\) summands, proves both assertions.

\section{The ratio tends to \texorpdfstring{\(\zeta(3)\)}{zeta(3)}}

\begin{proposition}\label{prop:ratio}
The closed forms \eqref{eq:ab} satisfy
\[
 \lim_{n\to\infty}\frac{b_n}{a_n}=\zetaThree.
\]
\end{proposition}

\begin{proof}
The ratio is a weighted average:
\begin{equation}\label{eq:average}
 \frac{b_n}{a_n}
 =
 \sum_{k=0}^np_{n,k}
 \bigl(2\HH n3-\HH k3\bigr).
\end{equation}
Fix a sufficiently small \(\delta>0\).  In the central region
\[
 \left|\frac{k}{n}-\frac1{\sqrt2}\right|<\delta
\]
we have \(k\geq cn\) for some \(c>0\).  Uniformly throughout this region,
\[
 \HH n3\longrightarrow\zetaThree,
 \qquad
 \HH k3\longrightarrow\zetaThree,
\]
and hence
\[
 2\HH n3-\HH k3\longrightarrow\zetaThree.
\]
The complementary region has total \(p_{n,k}\)-mass tending to zero by
\eqref{eq:concentration}, while all the harmonic expressions are uniformly
bounded.  Taking the weighted average in \eqref{eq:average} proves the claim.
\end{proof}

\begin{remark}
Writing \(T_m=\zetaThree-\HH m3\), the summandwise error is
\[
 \zetaThree-\bigl(2\HH n3-\HH k3\bigr)=2T_n-T_k.
\]
Since \(T_m=\frac1{2m^2}+O(m^{-3})\), its leading term vanishes exactly when
\(k\sim n/\sqrt2\), the maximum of the binomial kernel.  This explains the
visible first-order alignment.  The much stronger exponential cancellation
will come from the recurrence.
\end{remark}

\section{The denominator bound is termwise}

Let
\[
 d_n:=\lcm(1,2,\ldots,n).
\]
For every \(k\leq n\),
\[
 d_n^3\HH k3
 =
 \sum_{j=1}^k\frac{d_n^3}{j^3}
 =
 \sum_{j=1}^k\left(\frac{d_n}{j}\right)^3
 \in\mathbb Z.
\]
Because every \(K_{n,k}\) is an integer, the closed form immediately gives
\begin{equation}\label{eq:denominator}
 \boxed{d_n^3b_n\in\mathbb Z.}
\end{equation}
Of course \(a_n\in\mathbb Z\) as well.  This is the arithmetic advantage of
the short harmonic formula: no cancellation between different values of
\(k\) is needed to see the required denominator.

\section{The recurrence reveals the exponential cancellation}

The following calculation is the hinge of the proof.  Define the discrete
Wronskian
\[
 D_n:=b_na_{n-1}-b_{n-1}a_n.
\]
Solving \eqref{eq:recurrence} for \(u_{n+1}\) and applying it once to \(a_n\)
and once to \(b_n\) gives
\[
 D_{n+1}
 =\left(\frac{n}{n+1}\right)^3D_n.
\]
Since \(D_1=b_1a_0-b_0a_1=6\), induction yields the exact identity
\begin{equation}\label{eq:wronskian}
 \boxed{D_n=\frac6{n^3}.}
\end{equation}
It follows that
\begin{equation}\label{eq:increment}
 \frac{b_n}{a_n}-\frac{b_{n-1}}{a_{n-1}}
 =
 \frac6{n^3a_na_{n-1}}.
\end{equation}
Proposition~\ref{prop:ratio} identifies the limit of these increasing ratios.
Summing \eqref{eq:increment} from \(n+1\) to infinity therefore gives
\begin{equation}\label{eq:positive-tail}
 \boxed{
 \zetaThree-\frac{b_n}{a_n}
 =
 6\sum_{m=n+1}^{\infty}
 \frac1{m^3a_ma_{m-1}}.
 }
\end{equation}
In particular, the error is strictly positive.

Equation \eqref{eq:a-growth} shows that the \(m\)-th summand on the
right-hand side of \eqref{eq:positive-tail} has \(m\)-th root tending to
\(\lambda^{-2}\).  A positive tail of a sequence with exponential rate
\(\rho<1\) has the same exponential rate \(\rho\); hence
\[
 \lim_{n\to\infty}
 \left(\zetaThree-\frac{b_n}{a_n}\right)^{1/n}
 =\lambda^{-2}.
\]
For the linear form
\[
 L_n:=a_n\zetaThree-b_n
\]
we conclude that
\begin{equation}\label{eq:L-rate}
 \boxed{
 \lim_{n\to\infty}L_n^{1/n}
 =\lambda^{-1}
 =17-12\sqrt2
 \approx0.0294373.
 }
\end{equation}
Thus the recurrence has converted the signed cancellation inside the harmonic
sum into the manifestly positive tail \eqref{eq:positive-tail}.

\section{Irrationality}

\begin{theorem}
The number \(\zeta(3)\) is irrational.
\end{theorem}

\begin{proof}
The standard least-common-multiple estimate
\begin{equation}\label{eq:lcm-growth}
 \log d_n=n+o(n)
\end{equation}
is equivalent to the prime number theorem in its Chebyshev-function form.
Combining \eqref{eq:L-rate} and \eqref{eq:lcm-growth} gives
\[
 \lim_{n\to\infty}(d_n^3L_n)^{1/n}
 =
 \frac{e^3}{17+12\sqrt2}
 \approx0.591263<1.
\]
Consequently
\begin{equation}\label{eq:small}
 d_n^3L_n\longrightarrow0.
\end{equation}

Suppose, for a contradiction, that \(\zetaThree=p/q\) with integers
\(p,q>0\).  By \eqref{eq:denominator},
\[
 qd_n^3L_n
 =
 pd_n^3a_n-qd_n^3b_n
 \in\mathbb Z.
\]
The positive-tail formula \eqref{eq:positive-tail} says this integer is
strictly positive, whereas \eqref{eq:small} says it tends to zero.  For all
sufficiently large \(n\) it would therefore be an integer strictly between
\(0\) and \(1\), an impossibility.
\end{proof}

\begin{remark}
The full prime number theorem is more than is needed.  Any eventual estimate
\[
 d_n\leq C^n
 \qquad\text{with}\qquad
 C<(17+12\sqrt2)^{1/3}\approx3.23868
\]
is sufficient.
\end{remark}

\section{What the closed form contributes}

The irrationality mechanism is still Apéry's: a recurrence produces a large
integer solution and a very small companion solution.  The new formula does
not replace that mechanism.  Rather, it makes two of its least transparent
inputs almost effortless:
\[
 \frac{b_n}{a_n}\longrightarrow\zetaThree
 \quad\text{by concentration},\qquad
 d_n^3b_n\in\mathbb Z
 \quad\text{termwise}.
\]
Once those facts are visible, the discrete Wronskian completes the analytic
part in a few lines.  The result is a proof in which the arithmetic, the
limiting value, and the exponential cancellation each have a distinct and
simple explanation.

\end{document}
